\chapter{Системийн судалгаа}
\label{chapter1}
\section{Системийн тухай танилцуулга}
%Сонгож авсан системтэй холбоотой үйл ажиллагааны талаарх ерөнхий танилцуулгыг бичнэ. Яг өөрийн хөгжүүлэх системийн ажиллагааны тухай биш. Жишээлбэл, хувьцааны апп хөгжүүлэх бол энэ хэсэгт хувьцаа гэж юу болох, хувьцааг хэрхэн арилжаалдаг гэх мэт ерөнхий суурь ойлголтуудыг тайлбарлана. Харин яг өөрийнхөө хөгжүүлэх апп-тай холбоотой үйл ажиллагааг бүлэг 3-т бичнэ.

%Ашигласан материалыг references файлд хөтөлж явах ба ашиглаж буй хэсэгт cite команд ашиглан оруулна. Жишээ нь:
Ашиглаж буй материал \cite{internetusers} Энэ өгүүлбэрийг аль материалаас авч, ашиглаж буйг \cite{Reference1} ном зүй дэх агуулга заадаг \cite{ios}.

\section{Ижил төстэй системийн судалгаа}
\subsection{Монгол дахь ижил төстэй системүүд}
% Сонгож авсан сэдвийн хүрээнд дотоодод ажилладаг веб, апп байвал судалж, ерөнхий ажиллагаа, системийн үндсэн функц зэргийн тухай бичнэ. Мөн ажиллагааны ерөнхий агшныг харуулсан зураг оруулж болно. Яг сонгож авсан чиглэлд ажиллаж байгаа систем байхгүй бол төстэй систем судалж болно. Дор хаяж хоёр систем судалсан байх ёстой. 
\subsection{Гадаадын ижил төстэй системүүд}
% Гадаадад ижил төстэй систем олдох тул түүнтэй холбоотой судалсан мэдээллүүдийг оруулна. Дор хаяж хоёр систем судалсан байна.

%Судалсан системүүдийн талаар мэдээллийг нэгтгэсний дараа давуу ба сул талуудыг хүснэгтэд цэгцэлж харьцуулна. 
\section{Холбогдох хууль, дүрэм журам}
%системд хамааралтай хууль, дүрэм, заавар зэргийг судалж, шууд хамааралтай байж болох заалтуудыг оруулна.

\section{Шаардлага гаргаж авах техник}
\subsection{Хэрэглэгч, оролцогчдоос авсан бодит санал асуулга}
%Санал асуулгын асуултын тоо дор хаяж 8 байна.
%Санал асуулгыг дор хаяж 30 хүнээс авна.
%Санал асуулгын асуулт бүрийг энд оруулах ба асуулт дор хариултыг графикаар үзүүлнэ.
%Төгсгөлд нь санал асуулгаас гарсан үр дүнг нэгтгэж тайлбар бичнэ. 
\subsection{Ярилцлага}
%Системд хамааралтай оролцогчдоос ярилцлага авна. Ярилцлага авах хүний тоо 2.
%Ярилцлагын протокол хөтөлж, энд оруулна. Протоколд "хэзээ, хаана, хэнтэй ярилцаж буй, асуулт ба хариултын тэмдэглэлийг оруулна.


